% Zone „Das kennst du schon“ – aus bank/zufallsexperimente-und-pfadregeln/zone.jsonl (zusammenbau v0.1)
\einheitenkopf[zone][Kennst du schon]{Das kennst du schon}
\setcounter{aufgabe}{0}

%% TODO Titel: Ich-kann-Satz fehlt in der Bank (Platzhalter: Kettenname) – Nr. 1
%% TODO Anweisung: ein Satz über der ersten Teilaufgabe fehlt in der Bank – Nr. 1
\begin{aufgabe}{Ergebnismenge eines ein- oder zweistufigen Versuchs als geordnete Paare, Laplace-Regel einstufig, Summenregel, Gegenereignis, Zufallsgerät zu einer Wahrscheinlichkeit entwerfen}
\begin{teile}
\teil Eine Münze wird zweimal geworfen (K = Kopf, Z = Zahl). Alle Ergebnisse?
\teil Eine Urne mit 20 Kugeln soll P(rot) = $0{,}35$ und P(blau) = $0{,}25$ haben, der Rest ist gelb. Wie viele Kugeln je Farbe? rot \leerfeld, blau \leerfeld, gelb \leerfeld
\end{teile}
\end{aufgabe}

%% TODO Titel: Ich-kann-Satz fehlt in der Bank (Platzhalter: Kettenname) – Nr. 2
%% TODO Anweisung: ein Satz über der ersten Teilaufgabe fehlt in der Bank – Nr. 2
\begin{aufgabe}{Möglichkeiten vollständig aufzählen, Zählprinzip „mal“, Anordnungen einer kleinen Menge}
\begin{teile}
\teil Ein Imbiss hat 4 Brötchen und 3 Soßen. Wie viele Kombinationen aus einem Brötchen und einer Soße? \leerfeld
\teil Vier Freunde stellen sich für ein Foto in eine Reihe. Wie viele Reihenfolgen gibt es? \leerfeld
\end{teile}
\end{aufgabe}

%% TODO Titel: Ich-kann-Satz fehlt in der Bank (Platzhalter: Kettenname) – Nr. 3
%% TODO Anweisung: ein Satz über der ersten Teilaufgabe fehlt in der Bank – Nr. 3
\begin{aufgabe}{Brüche multiplizieren (auch mehrere Faktoren mit fallenden Nennern), kürzen, gleichnamig und ungleichnamig addieren; Bruch ↔ Dezimalzahl ↔ Prozent; Potenzen kleiner Brüche}
\begin{teile}
\teil $\frac{1}{4} \cdot \frac{2}{3}$ \leerfeld
\teil $\frac{5}{8} \cdot \frac{4}{7} \cdot \frac{3}{6}$ \leerfeld
\end{teile}
\end{aufgabe}

%% TODO Titel: Ich-kann-Satz fehlt in der Bank (Platzhalter: Kettenname) – Nr. 4
%% TODO Anweisung: ein Satz über der ersten Teilaufgabe fehlt in der Bank – Nr. 4
\begin{aufgabe}{Symbolschreibweise der Ereignisse (A ∩ B, A ∪ B, Ā, A \ B, $P_A(B))$ lesen; die bedingte Wahrscheinlichkeit als Quotient}
\begin{teile}
\teil Ein Würfel wird geworfen, A: „Die Augenzahl ist gerade“. Was bedeutet $\bar{A}$ (nicht A)?
\teil $P(A \cap B) = 0{,}12$ und $P(A) = 0{,}4$. $P_A(B)$ (B unter der Bedingung A)? $P_A(B) = $ \leerfeld
\end{teile}
\end{aufgabe}

%% TODO Titel: Ich-kann-Satz fehlt in der Bank (Platzhalter: Kettenname) – Nr. 5
%% TODO Anweisung: ein Satz über der ersten Teilaufgabe fehlt in der Bank – Nr. 5
\begin{aufgabe}{Anteil mal Gesamtzahl (absolute Anzahl), Mittelpunktswinkel als Anteil des Vollwinkels, Prozentsatz als Dezimalzahl}
\begin{teile}
\teil 20\,\% von 50 Losen sind Gewinne. Wie viele Gewinnlose? \leerfeld
\teil Ein Sektor soll mit Wahrscheinlichkeit $0{,}15$ getroffen werden. Mittelpunktswinkel? \leerfeld[°]
\end{teile}
\end{aufgabe}

%% TODO Titel: Ich-kann-Satz fehlt in der Bank (Platzhalter: Kettenname) – Nr. 6
%% TODO Anweisung: ein Satz über der ersten Teilaufgabe fehlt in der Bank – Nr. 6
\begin{aufgabe}{Baumdiagramm mit Zurücklegen oder mit zwei Geräten zeichnen und ergänzen (Äste an einem Punkt ergeben eins), Pfadregel, Summenregel, „mindestens einmal“ über das Gegenereignis, dreistufig}
\begin{teile}
\teil Urne mit 2 roten und 3 blauen Kugeln, zweimal mit Zurücklegen. P(rot, rot)? P = \leerfeld
\teil Ein Glücksrad (A oder B) wird zweimal gedreht. Ergänze die fehlenden Äste.

\baumzwei{A/0{,}3,B/}{A/,B/}{A/,B/}
\end{teile}
\end{aufgabe}

%% TODO Titel: Ich-kann-Satz fehlt in der Bank (Platzhalter: Kettenname) – Nr. 7
%% TODO Anweisung: ein Satz über der ersten Teilaufgabe fehlt in der Bank – Nr. 7
\begin{aufgabe}{Ziehen ohne Zurücklegen zweistufig (Nenner um eins kleiner, bei der gezogenen Sorte auch der Zähler), zweite Person zieht, dreistufig ein Pfad, „unter den ersten drei“}
\begin{teile}
\teil Urne mit 5 roten und 3 blauen Kugeln, zweimal ohne Zurücklegen. P(rot, rot)? P = \leerfeld
\teil 10 Kugeln, davon 4 gelb. Drei werden ohne Zurücklegen gezogen. P(dreimal gelb)? P = \leerfeld
\end{teile}
\end{aufgabe}

%% TODO Titel: Ich-kann-Satz fehlt in der Bank (Platzhalter: Kettenname) – Nr. 8
%% TODO Anweisung: ein Satz über der ersten Teilaufgabe fehlt in der Bank – Nr. 8
\begin{aufgabe}{Kombinatorik}
\begin{teile}
\teil $6!$ (6 Fakultät)? \leerfeld
\teil Aus 8 Personen werden 3 für ein Team ausgewählt. Wie viele Teams sind möglich? \leerfeld
\end{teile}
\end{aufgabe}

%% TODO Titel: Ich-kann-Satz fehlt in der Bank (Platzhalter: Kettenname) – Nr. 9
%% TODO Anweisung: ein Satz über der ersten Teilaufgabe fehlt in der Bank – Nr. 9
\begin{aufgabe}{Lineare Gleichungen, quadratische Gleichungen (Lösungsformel, Lösungen sieben) und einfache Bruchgleichungen aufstellen und lösen}
\begin{gleichungsraster}[2]
\gl{0{,}5x + 0{,}2 = 0{,}6} &
\gl{x^2 - x + 0{,}24 = 0} \\
\end{gleichungsraster}
\end{aufgabe}

%% TODO Titel: Ich-kann-Satz fehlt in der Bank (Platzhalter: Kettenname) – Nr. 10
%% TODO Anweisung: ein Satz über der ersten Teilaufgabe fehlt in der Bank – Nr. 10
\begin{aufgabe}{Ziehen ohne Zurücklegen zweistufig (Nenner um eins kleiner, bei der gezogenen Sorte auch der Zähler), zweite Person zieht, dreistufig ein Pfad, „unter den ersten drei“ – Fehler finden}
\begin{teile}
\teil Urne mit 6 roten und 4 grünen Kugeln, zweimal ohne Zurücklegen. Lisa rechnet P(grün, grün) $= \frac{4}{10} \cdot \frac{4}{10} = 0{,}16$. Finde den Fehler.
\end{teile}
\end{aufgabe}

%% TODO Titel: Ich-kann-Satz fehlt in der Bank (Platzhalter: Kettenname) – Nr. 11
%% TODO Anweisung: ein Satz über der ersten Teilaufgabe fehlt in der Bank – Nr. 11
\begin{aufgabe}{Ziehen ohne Zurücklegen zweistufig (Nenner um eins kleiner, bei der gezogenen Sorte auch der Zähler), zweite Person zieht, dreistufig ein Pfad, „unter den ersten drei“ – selbst rechnen}
\begin{teile}
\teil Urne mit 5 weißen und 3 schwarzen Kugeln, zweimal ohne Zurücklegen. P(schwarz, schwarz)? P = \leerfeld
\end{teile}
\end{aufgabe}
